\documentclass[12pt]{article}
\usepackage[T1]{fontenc}
\usepackage[utf8]{inputenc}
\usepackage{lmodern}
\usepackage{textcomp}
\usepackage{enumitem}
\usepackage{geometry}
\geometry{margin=1in}
\begin{document}
\begin{center}
Answer Key - Political Science - Graduate Level Assessment 12 \
\small (Answers are ordered exactly as in the question paper.)
\end{center}

\section*{Multiple Choice}
\begin{enumerate}[label=\arabic*.]
\item \textbf{Answer:} A
\\ \textit{Options:} A) Critical questioning.; B) Realist logic.; C) State-centric security analysis.; D) All of these options.
\\ \textit{Note:} Critical questioning represents a move away from traditional interpretations of security studies by challenging established norms and assumptions, thereby broadening the understanding of security beyond state-centric and military-focused perspectives.
\item \textbf{Answer:} B
\\ \textit{Options:} A) Jimmy Carter; B) Ronald Reagan; C) Bill Clinton; D) Dwight Eisenhower
\\ \textit{Note:} The correct answer is Ronald Reagan because his administration implemented a significant military buildup during the 1980 s, resulting in the largest peacetime defense spending increase in U.S. history as part of his Cold War strategy to counter Soviet influence.
\item \textbf{Answer:} C
\\ \textit{Options:} A) It feared the League would encourage Soviet influence in the US; B) It feared the League would be anti-democratic; C) It feared the League would commit the US to an international alliance; D) Both a and b
\\ \textit{Note:} Congress opposed Wilson's proposal for the League of Nations because it was concerned that joining the League would obligate the United States to participate in international conflicts and entanglements, undermining its sovereignty and the principle of isolationism that many lawmakers favored at the time.
\item \textbf{Answer:} A
\\ \textit{Options:} A) Members of Congress negotiating directly with foreign governments; B) Face-to-face meetings between state leaders; C) The president consulting Congress on foreign policy issues; D) Bilateral talks that do not involve a third-party negotiator
\\ \textit{Note:} Direct diplomacy refers to the practice where members of Congress engage in direct negotiations with foreign governments, thereby bypassing traditional diplomatic channels typically reserved for the executive branch, highlighting the legislative branch's role in foreign policy.
\item \textbf{Answer:} D
\\ \textit{Options:} A) As a means of exerting economic power to advance peace and freedom.; B) To end economic support for morally unacceptable regimes.; C) To isolate dangerous regimes and weaken threats to regional or global security.; D) All of these options.
\\ \textit{Note:} The correct answer "All of these options" is valid because sanctions serve multiple functions, including influencing behavior, signaling disapproval, and exerting economic pressure, all of which contribute to their effectiveness in international relations.
\end{enumerate}

\section*{True / False}
\begin{enumerate}[label=\arabic*.]
\setcounter{enumi}{5}
\item \textbf{Answer:} False
\\ \textit{Options:} A) True; B) False
\\ \textit{Note:} Human society is characterized by constant change driven by factors such as cultural evolution, social movements, technological advancements, and political shifts, which create unpredictable dynamics rather than inherent stability.
\item \textbf{Answer:} True
\\ \textit{Options:} A) True; B) False
\\ \textit{Note:} Democratic enlargement encompasses the strategy of promoting democracy and free market principles worldwide, reflecting the belief that these systems contribute to global stability and prosperity.
\item \textbf{Answer:} True
\\ \textit{Options:} A) True; B) False
\\ \textit{Note:} The responsibility to prevent is deemed the most crucial pillar of the Responsibility to Protect (R$_{2}$P) framework because it emphasizes the proactive measures that states and the international community must take to avert mass atrocities before they occur, thereby addressing the root causes of conflict and protecting populations from harm.
\item \textbf{Answer:} True
\\ \textit{Options:} A) True; B) False
\\ \textit{Note:} Political policies that alter demographic characteristics, the criminalization of cultural practices, and military force for ethnic cleansing undermine social cohesion, violate human rights, and threaten the stability and diversity of societies, making them significant threats to social order and peace.
\item \textbf{Answer:} True
\\ \textit{Options:} A) True; B) False
\\ \textit{Note:} John Foster Dulles coined the term 'brinkmanship' to describe the policy of pushing a dangerous situation to the edge of disaster in order to achieve favorable outcomes in the Cold War, particularly in relation to the Soviet Union.
\end{enumerate}

\section*{Long Form Response}
\begin{enumerate}[label=\arabic*.]
\setcounter{enumi}{10}
\item \textbf{Answer:} The end of the Cold War marked a significant turning point in the re-evaluation of security policies, particularly through the lens of environmental security and multilateralism. As the bipolar tension diminished, states began to recognize that traditional security threats were not solely military but included environmental degradation and resource scarcity, which posed risks to global stability. This shift fostered an increased emphasis on multilateral cooperation, as countries acknowledged that collective action was essential to address transnational environmental challenges. Consequently, security studies expanded to incorporate frameworks that integrate ecological considerations, reflecting a broader understanding of security that encompasses human survival and well-being. Ultimately, the re-conceptualization of security post-Cold War underscored the interdependence of environmental sustainability and international peace, facilitating collaborative efforts to tackle pressing global issues.
\\ \textit{Note:} The answer correctly identifies that the end of the Cold War prompted a broader understanding of security that included environmental issues and resource management, leading to enhanced multilateral cooperation among states to address these non-traditional threats to global stability.
\item \textbf{Answer:} In political science, the concept of 'nations' is often defined by shared characteristics such as language, culture, ethnicity, and historical experiences, which contribute to a collective identity. While nations may aspire to self-determination and sovereignty, their relationship with the boundaries and borders of the state is complex and multifaceted. In some cases, a nation may coincide with state boundaries, leading to a nation-state where the political and cultural identities align. However, there are numerous instances where nations exist across state borders or where multiple nations inhabit a single state, complicating the notion of sovereignty and governance. Thus, while nations are defined by cultural and social constructs, their relationship with state boundaries can reveal tensions and challenges in political organization and identity formation.
\\ \textit{Note:} The answer correctly identifies that nations are defined by shared characteristics contributing to a collective identity, and it highlights the complex relationship between nations and state boundaries, emphasizing that while they can align in a nation-state, this is not always the case.
\end{enumerate}

\vfill
\noindent\textit{End of Answer Key}
\end{document}